\documentclass[sigconf]{acmart}

\usepackage{graphicx}
\usepackage{balance}
\usepackage{amsmath}
\usepackage{amssymb}
\usepackage{booktabs}

\setcopyright{acmcopyright}
\copyrightyear{2026}
\acmYear{2026}
\acmConference[MM '26]{Proceedings of the 34th ACM International Conference on Multimedia}{October 2026}{Rio de Janeiro, Brazil}

\begin{document}

\title{EDRA: Experience-Driven Rule Adaptation for Self-Improving Outreach Agents}

\author{Daniil Onishchenko}
\email{danial92335@mail.ru}
\affiliation{%
   \institution{Kazan Federal University}
   \city{Kazan}
   \country{Russia}}

\author{Aleksandr Farseev}
\email{farseev@itmo.ru}
\affiliation{%
  \institution{ITMO University}
  \city{Saint Petersburg}
  \country{Russia}
}
\affiliation{%
  \institution{SoMin.ai Research}
  \country{Singapore}
}

\begin{abstract}
LLM-based outreach agents personalize messages at scale, yet their behavioral strategies---framing, tone, level of assertiveness---remain static prompt templates while the audiences they address are heterogeneous and shifting.
Recent systems (ExpeL, EvolveR, ReasoningBank) close parts of this feedback loop by extracting task-general strategies from interaction histories, but none conditions rules on \textit{who} the agent is talking to: their insights are global, not audience-aware.
We present EDRA (Experience-Driven Rule Adaptation), a closed-loop architecture that introduces \textit{cluster-conditional} rule induction.
Visitor profiles are embedded via MiniLM, clustered with HDBSCAN, and matched to behavioral rules through weighted KNN voting; an LLM revises rules when cluster-level consistency degrades.
We frame rule application as a contextual bandit and validate offline on 976 real cold-outreach interactions from a commercial CRM via a doubly robust estimator.
EDRA's cluster-conditional policy yields an estimated reward of \textbf{[X.XXX]}, outperforming the best single-strategy baseline (\textbf{[X.XXX]}), confirming that optimal outreach strategy differs by audience segment.
At the conference booth, each attendee simultaneously experiences and trains the system.
\end{abstract}

\keywords{LLM agents, rule induction, contextual bandits, profile clustering, off-policy evaluation, interactive demonstration}

\maketitle

\begin{figure*}[t]
    \centering
    \includegraphics[width=0.9\textwidth]{EDRA_workflow.png}
    \caption{EDRA architecture. Profiles are embedded via MiniLM into $\mathbb{R}^{384}$. HDBSCAN discovers clusters; KNN votes across the $K$ nearest profiles to select a behavioral rule. Episodes feed into consistency monitoring, which triggers LLM revision when performance degrades. The same loop processes booth dialogues and asynchronous outreach.}
    \label{fig:architecture}
\end{figure*}

% ============================================================
\section{Introduction}
% ============================================================

The behavioral strategies that LLM-based outreach agents follow---framing, tone, how hard to push---are hand-written prompt templates.
The audiences they address, however, are neither uniform nor stable.
A ``collaborative exploration'' pitch that resonates with early-career researchers reads as naive to a VP of Engineering; the agent treats both identically.
This is the core asymmetry: the agent personalizes \textit{content} through an LLM, but the \textit{strategy} that governs content generation remains one-size-fits-all.

Recent work has begun closing parts of this loop.
TRAD~\cite{trad2024} retrieves past reasoning steps but does not abstract them into reusable rules.
MetaFlowLLM~\cite{metaflowllm2025} induces hybrid workflows offline, without online revision.
ExpeL~\cite{expel2024} extracts task-general ``insights'' from episodes in an offline training phase with no degradation monitoring.
EvolveR~\cite{evolver2025} and ReasoningBank~\cite{reasoningbank2025} close the feedback loop with online revision, but their rules are task-general: a single policy for all users, not conditioned on who the agent is talking to.
The gap is specific: no existing system induces behavioral rules \textit{per audience segment} and revises them when segment-level performance degrades.

We present \textbf{EDRA} (Experience-Driven Rule Adaptation).
First, we introduce an architecture where behavioral rules are induced per profile cluster rather than globally, so distinct audience types receive distinct strategies.
Building on this, we formalize a KNN-based rule application mechanism where new visitors are matched to cluster rules through weighted cosine voting over semantically similar profiles.
Finally, we validate through offline evaluation on 976 real outreach interactions framed as a contextual bandit~\cite{li2010contextual}, comparing four policies via a doubly robust estimator~\cite{dudik2011doubly}, and through live booth interactions where each attendee simultaneously experiences and trains the system.

% ============================================================
\section{The EDRA Architecture}
% ============================================================

\noindent\textbf{Profile Embedding.}
Figure~\ref{fig:architecture} illustrates the full pipeline.
Let $\mathcal{P} = \{p_1, \ldots, p_N\}$ denote accumulated visitor profiles.
A deterministic summarization $\text{summ}(\cdot)$ concatenates profile fields, and a sentence model $\phi$ (all-MiniLM-L6-v2) maps the result to a unit vector:
\begin{equation}
    \mathbf{v}_i = \frac{\phi(\text{summ}(p_i))}{\|\phi(\text{summ}(p_i))\|_2} \in \mathbb{R}^{384}
    \label{eq:embedding}
\end{equation}
We cluster profiles, not episodes: who a person \textit{is} determines cluster membership; how they \textit{responded} determines whether the cluster's rule works.

\noindent\textbf{Clustering and Rule Induction.}
HDBSCAN~\cite{campello2013density} discovers density-based groups over the embedding space without a fixed cluster count.
When a cluster accumulates enough episodes, the LLM receives representative interactions and generates a five-slot behavioral rule $\rho_k$: \textit{framing}, \textit{tone}, \textit{opener\_type}, \textit{word\_target}, \textit{ask\_size}.
Slots are static (literal values) or dynamic (LLM sub-prompts filled at application time).

\noindent\textbf{KNN Rule Application.}
Hard cluster assignment is brittle: boundary profiles get assigned arbitrarily and noise points receive no rule.
Given a new visitor $p_{\text{new}}$ with embedding $\mathbf{v}_{\text{new}}$, let $\mathcal{N}_K(p_{\text{new}})$ denote the $K$ nearest profiles by cosine similarity. The weighted vote for cluster $C_k$:
\begin{equation}
    w_k = \sum_{\substack{p_{(j)} \in \mathcal{N}_K(p_{\text{new}}) \\ p_{(j)} \in C_k}} \mathbf{v}_{\text{new}}^\top \mathbf{v}_{(j)}
    \label{eq:knn_vote}
\end{equation}
The winning cluster's rule is applied: $\rho^* = \rho_{\arg\max_k w_k}$.
If no neighbor has an active rule, the agent improvises ($K{=}7$).

\noindent\textbf{Consistency Monitoring and Revision.}
For rule $\rho_k$ induced at time $\tau_{\text{ind}}$, the consistency score over post-induction episodes:
\begin{equation}
    \text{CS}(\rho_k, t) = \frac{|\{e \in \mathcal{E}_{C_k} : \tau_e > \tau_{\text{ind}},\; o_e = \texttt{accepted}\}|}{|\{e \in \mathcal{E}_{C_k} : \tau_e > \tau_{\text{ind}}\}|}
    \label{eq:cs}
\end{equation}
When CS drops below $\theta_{\text{revise}}$, the LLM receives the rule and contradicting episodes, produces a revision, and streams its reasoning via Server-Sent Events.

% ============================================================
\section{Offline Validation}
% ============================================================

We frame EDRA's rule application as a contextual bandit: the context $x$ is a profile embedding, the action $a \in \{\textit{cold\_template}, \textit{cold\_personal}, \textit{feature\_ann.}\}$ is the outreach strategy, and the reward $r \in \{0, 1\}$ is whether the recipient replied.
We evaluate on 976 first-contact interactions extracted from a commercial CRM, filtered to cold outreach only. Propensity scores $\pi_0(a \mid x)$ are reconstructed from the campaign assignment patterns in the logged data.

\noindent\textbf{Level 1: Clustering Quality.}
HDBSCAN over MiniLM embeddings of enriched profiles yields \textbf{[N]} clusters with a silhouette score of \textbf{[X.XXX]}.
A $\chi^2$ test confirms that reply rates differ significantly across clusters ($p < $ \textbf{[X.XXX]}), establishing that audience heterogeneity is real, not an artifact of the embedding.

% ---- PLACEHOLDER: insert silhouette and chi-squared results ----

\noindent\textbf{Level 2: Does Optimal Strategy Differ by Cluster?}
This is the central test. We compare four policies via the doubly robust estimator~\cite{dudik2011doubly}:
$\pi_{\text{uniform}}$ (always \textit{cold\_template}),
$\pi_{\text{random}}$,
$\pi_{\text{best\_single}}$ (the single strategy with the highest overall reply rate),
and $\pi_{\text{edra}}$ (cluster-conditional).
If $\pi_{\text{edra}}$ cannot beat $\pi_{\text{best\_single}}$, cluster-conditioning adds nothing.
The estimated rewards are
\textbf{[X.XXX]} ($\pi_{\text{uniform}}$),
\textbf{[X.XXX]} ($\pi_{\text{random}}$),
\textbf{[X.XXX]} ($\pi_{\text{best\_single}}$),
and \textbf{[X.XXX]} ($\pi_{\text{edra}}$).

% ---- PLACEHOLDER: insert DR estimator results for 4 policies ----

\noindent\textbf{Level 3: Learning Curve.}
Following the prequential protocol~\cite{gama2013evaluating}, we sort episodes chronologically and evaluate in expanding batches, comparing EDRA against LinUCB~\cite{li2010contextual}.
\textbf{[TBD: whether the estimated reward improves monotonically with batch size and how EDRA compares to the LinUCB baseline.]}

% ---- PLACEHOLDER: insert prequential learning curve description ----

% ============================================================
\section{Becoming the Subject}
% ============================================================

\begin{figure}[t]
    \centering
    \includegraphics[width=\columnwidth]{edra_UI.png}
    \caption{The EDRA booth. An anime-style agent conducts a personalized visual-novel dialogue. The Expert View reveals the rulebook, cluster scatter plot, and streaming revision console.}
    \label{fig:booth}
\end{figure}

The booth makes the closed loop tangible (Figure~\ref{fig:booth}).
An attendee enters their LinkedIn URL; the system fetches and embeds their profile, performs KNN lookup, and applies the matched rule---or improvises if none exists.
The dialogue proceeds over 3--7 turns via three contextual response buttons generated by the LLM, with an inner-monologue mechanic that reveals the agent's strategic reasoning between turns.
A 12-state emotion system drives avatar expressions through the conversation.
After the session, the episode is persisted and the profile is purged within one hour, retaining only the anonymized embedding.

The attendee is simultaneously a user and a data point: their interaction trains the rule base that will shape future conversations.
The Expert View panel exposes induced rules, a t-SNE cluster scatter plot with KNN neighbors highlighted, and a streaming reflection console showing revision reasoning in real time.
The system runs as a single-process FastAPI application with three asyncio loops, requiring only a laptop.

% ============================================================
\section{Conclusion}
% ============================================================

EDRA closes the feedback loop between episodic experience and behavioral strategy, conditioning rules on audience type rather than treating all visitors identically.
The architecture is channel-agnostic: the same induction--application--revision cycle processes both live booth dialogues and asynchronous email outreach.
Closing this loop is, in our view, a direction worth pursuing well beyond outreach---wherever an LLM agent interacts with heterogeneous audiences.
Looking ahead, two directions shape future work: (1)~grounding cluster labels in empirical audience data for interpretable segment descriptions, and (2)~extending the action space from single-message strategies to multi-turn dialogue policies.

\balance
\bibliographystyle{ACM-Reference-Format}
\bibliography{edra_demo}

\end{document}
